% !TEX program = xelatex
% AI-effects 프로젝트 --- Sensor Tower True Audience 이용자 수로 만든 제품별·국가별·Worldwide 월간 메시지 수,
% 업무용 메시지 수, 업무용 대화량 데이터의 원자료와 구축 과정.
% 표 본문: results/table/17_{anchor,products,country,worldwide}.tex
%          (code/04_analysis_sensortower_messages_summary.do가 생성)
% 데이터:  code/02_clean_sensortower_share_all_products.do, code/03_merge_sensortower_openai_messages.do
% 컴파일: xelatex -> biber -> xelatex -> xelatex (kotex, biblatex/biber 필요)
\documentclass[11pt]{article}

\usepackage{fontspec}
\usepackage{kotex}
\usepackage[a4paper,margin=2.2cm]{geometry}
\usepackage{longtable}
\usepackage{booktabs}
\usepackage{array}
\usepackage{amsmath}
\usepackage{amssymb}
\usepackage{xcolor}
\usepackage{hyperref}
\usepackage{enumitem}
\usepackage{titlesec}
\usepackage[backend=biber,style=authoryear,maxcitenames=2,uniquelist=false,sorting=nyt]{biblatex}
\addbibresource{../../reference/references.bib}

\hypersetup{colorlinks=true,linkcolor=blue!50!black,citecolor=blue!50!black,urlcolor=blue!50!black}
\newcolumntype{L}[1]{>{\raggedright\arraybackslash}p{#1}}
\newcolumntype{R}[1]{>{\raggedleft\arraybackslash}p{#1}}
\setlength{\LTleft}{0pt}
\setlength{\LTright}{0pt}
\renewcommand{\arraystretch}{1.15}
\setlength{\tabcolsep}{4pt}

\title{Sensor Tower 이용자 수 기반 AI 어시스턴트 월간 메시지$\cdot$업무용 대화 데이터 구축\\[2mm]
\large --- 원자료와 구축 과정 ---}
\author{AI-effects 프로젝트}
\date{\today}

\begin{document}
\maketitle

\begin{abstract}
\noindent 본 문서는 Sensor Tower \textit{State of AI 2026}\parencite{sensortower-2026-state-of-ai}의 True Audience 이용자 수, \textcite{chatterji-2025-how-people-use-chatgpt}의 ChatGPT 주간 메시지 수, OpenAI Signals v2.0\parencite{openai-2026-signals-v2}의 국가별 업무용 메시지 비중을 결합해 만든 다섯 가지 데이터---(1) 제품별 월간 메시지 수, (2) 국가별 월간 메시지 수, (3) Worldwide 월간 메시지 수, (4) 국가별 월간 업무용 메시지 수, (5) Worldwide 월간 업무용 메시지 수와 업무용 대화량---의 원자료와 구축 과정을 기록한다. 구축은 두 단계다. 1단계(\texttt{02\_clean\_sensortower\_share\_all\_products.do})는 True Audience 이용자 수를 시장$\times$월로 합산해 국가별 점유율을 만들고, Worldwide 안의 제품별 점유율을 다시 계산해 Sensor Tower 공개 점유율과 대조한다. 두 값은 304개 셀 모두에서 0.11\%p 안에서 일치한다. 2단계(\texttt{03\_merge\_sensortower\_openai\_messages.do})는 ``모든 제품의 이용자 1인당 메시지 수가 같다''와 ``모든 제품의 업무용 메시지 비중이 같다''는 두 가정 아래, ChatGPT 2025년 7월 값에서 얻은 1인당 월 메시지 수(약 80건)를 이용자 수에 곱하고 국가별 ChatGPT 업무 비중을 적용한다. 그 결과 8개 제품 합계 월간 메시지는 2025년 7월 1,374억 건, 2026년 5월 1,916억 건이고, 2026년 5월 업무용 메시지는 612억 건(31.9\%), 대화당 메시지 3건이면 업무용 대화 204억 건이다. 이 값들은 두 가정에 전적으로 의존하므로 ``관측값''이 아니라 ``가정 아래의 규모 추정''으로 읽어야 한다.
\end{abstract}

\tableofcontents

% ============================================================
\section{목적과 개요}
% ============================================================

이 저장소의 목표는 AI 노출도(exposure)에 실제 사용량(usage)을 결합하는 것이다. 사용량 자료는 대부분 \textbf{비중}(점유율, 업무 비중, 과업 구성)만 공개하고 \textbf{규모}(건수)는 공개하지 않는다. 본 문서의 데이터는 이 빈칸을 메우기 위해, 가장 신뢰할 만한 한 시점의 규모값(ChatGPT 2025년 7월 주 메시지 180억 건)을 Sensor Tower의 제품별$\cdot$국가별 이용자 수로 넓혀 \textbf{제품$\times$국가$\times$월 단위의 메시지 규모}를 만든다.

구축 흐름은 다음과 같다.
\begin{enumerate}[leftmargin=*]
  \item \textbf{점유율}(1단계): True Audience 이용자 수 $\to$ 국가별 전체 제품 점유율(⑮), 제품별 Worldwide 점유율(⑯).
  \item \textbf{1인당 메시지}: ChatGPT 2025년 7월 월 메시지 $\div$ ChatGPT 2025년 7월 이용자.
  \item \textbf{메시지 수}: 1인당 메시지 $\times$ 이용자 수 $\to$ 제품별, 국가별, Worldwide.
  \item \textbf{업무용 메시지}: 국가 메시지 $\times$ 그 국가의 ChatGPT 업무 비중(Signals).
  \item \textbf{대화량}: 업무용 메시지 $\div$ 대화당 메시지 수 $k \in \{1.5, 3, 5, 7, 10\}$.
\end{enumerate}
번호 ①--⑯은 원자료 폴더의 \texttt{\_download\_manifest.md}와 report 13의 데이터셋 번호를 따른다.

% ============================================================
\section{원자료}
% ============================================================

\begin{longtable}{L{2.6cm}L{4.6cm}L{4.4cm}L{4.2cm}}
\caption{원자료}\label{tab:raw}\\
\toprule
자료 & 파일 & 내용 & 단위$\cdot$범위 \\
\midrule
\endfirsthead
\toprule
자료 & 파일 & 내용 & 단위$\cdot$범위 \\
\midrule
\endhead
Sensor Tower ① & \texttt{sensortower\_true\_audience\_monthly\_long.csv} & 제품별 중복 제거 독립 앱+웹 월간 이용자 수(True Audience) & 명. 19개 시장(Worldwide + 18개국) $\times$ 8개 제품 $\times$ 2023-04--2026-05 \\
Sensor Tower ⑥ & \texttt{sensortower\_true\_audience\_share\_monthly\_long.csv} & ①의 제품별 점유율(Sensor Tower 공개값) & \%, 소수 첫째 자리. 검증용 \\
Chatterji et al.\ (2025) & \texttt{scaling\_parameters\_public.csv}의 \texttt{chatterji\_msgs\_week\_2025\_07} & ChatGPT 주 메시지 180억 건(2025년 7월, p.1) & 백만 건 \\
OpenAI Signals v2.0 & \texttt{share\_of\_messages\_by\_work\_related\_country\_month.csv}, \texttt{\ldots\_work\_related\_month.csv} & ChatGPT 메시지 중 업무 관련(1)/비업무(0) 비중 & 비율(0--1). 국가별$\cdot$전 세계, 2024-07--2026-06 \\
\bottomrule
\end{longtable}

\subsection{Sensor Tower True Audience(①, ⑥)}
Sensor Tower \textit{State of AI 2026}의 차트 ``Deduplicated Standalone App \& Web Users for Top AI Assistants''에서 추출한 값이다(report 13). 한 제품 안에서 앱과 웹 이용자의 중복을 제거한 월간 이용자 수이며, 8개 항목(ChatGPT, Google Gemini, Claude, Microsoft Copilot, Perplexity, DeepSeek, Grok, Other)으로 나뉜다. 다른 앱에 내장된 AI 이용자(예: WhatsApp 안의 Meta AI)는 빠진다.
\begin{itemize}
  \item \textbf{Worldwide는 25개 시장 합계}이고 국가 시트는 18개국뿐이다. 나머지 7개 시장은 이름이 공개되지 않아 ``Worldwide $-$ 18개국 합계''로 한 행(residual)을 만든다.
  \item 여러 어시스턴트를 쓰는 사람은 \textbf{제품마다 한 번씩} 센다. 제품을 합산한 값은 사람 수가 아니라 ``제품별 이용자 수의 합''이다.
  \item 값은 유효숫자 약 3자리로 반올림되어 있다. ⑥(공개 점유율)은 소수 첫째 자리 값이다.
\end{itemize}

\subsection{Chatterji et al.\ (2025)의 기준값}
OpenAI 연구진의 NBER 논문 \textcite{chatterji-2025-how-people-use-chatgpt}는 ``2025년 7월까지 7억 명이 매주 180억 건의 메시지를 보냈다''고 쓴다(p.1). 범위는 소비자 요금제(Free, Plus, Pro)이며, 이 값의 선택 이유와 한계는 report 15에서 다루었다. 월간 환산은 report 15와 같이 ``주간 $\times$ 그 달 일수 $\div$ 7''로 한다.

\subsection{OpenAI Signals v2.0 업무 관련 비중}
Signals는 매월 소비자 ChatGPT 메시지 표본(약 30만 건)을 LLM 분류기로 분류해 점유율만 공개한다(report 12). \texttt{work\_related} = 1은 업무 관련 메시지다. 국가별 파일은 135개국을 포함하며 Sensor Tower 18개국은 24개월 모두 있다. 국가 코드는 ISO 3166-1 alpha-2, 월은 \texttt{YYYY-MM-01} 형식이다. 차등정보보호 잡음이 들어 있다.

% ============================================================
\section{1단계: 점유율 데이터(\texttt{02\_clean\_sensortower\_share\_all\_products.do})}
% ============================================================

\subsection{국가별 전체 제품 점유율(⑮)}
월 $t$, 시장 $c$에 대해 8개 제품 이용자 수를 합한다.
\[
  \text{all\_users}_{c,t} = \sum_{p} \text{users}_{c,p,t}, \qquad
  \text{share\_all\_products\_pct}_{c,t} = 100 \times \frac{\text{all\_users}_{c,t}}{\text{all\_users}_{\text{WW},t}} .
\]
18개국에 residual 행을 더하면 합이 100\%가 된다(do 파일에서 검사). 보조 변수 \texttt{share\_within\_18\_pct}는 분모를 18개국 합계로 둔 값이다. 산출: \texttt{sensortower\_share\_all\_products\_monthly\_\{long.csv, long.dta, wide.csv\}}(긴 형식 722행 = 38개월 $\times$ 19행).

\subsection{제품별 Worldwide 점유율(⑯)과 공개값 대조}
Worldwide 시장 안에서 제품 $p$의 비중을 ①에서 다시 계산한다.
\[
  \text{share\_ww\_pct}_{p,t} = 100 \times \frac{\text{users}_{\text{WW},p,t}}{\sum_{q} \text{users}_{\text{WW},q,t}} .
\]
⑥의 Worldwide 행과 정의가 같으므로 두 값을 \texttt{assistant}$\times$\texttt{month}로 붙여 대조했다. 304개 셀(8개 제품 $\times$ 38개월)이 모두 짝지어졌고, 상관은 1.0000, 차이의 절댓값은 평균 0.025\%p, 최대 0.108\%p(ChatGPT 2026-02: 51.19\% 대 51.3\%)이다. ⑥의 반올림 폭(0.05\%p)을 넘는 24개 셀은 ①의 유효숫자 3자리 반올림으로 설명된다. 따라서 Sensor Tower의 공개 점유율은 ``제품 이용자 수 $\div$ 8개 항목 합계''로 정확히 재현되며, ①을 규모 계산의 바탕으로 써도 공개 점유율과 어긋나지 않는다. 산출: \texttt{sensortower\_share\_worldwide\_by\_product\_monthly\_long.\{csv, dta\}}(304행; 공개값 \texttt{share\_pct\_st}와 차이 \texttt{diff\_pp} 포함).

두 산출물은 사용자 지정에 따라 원자료 폴더(\texttt{data/raw/macro/scaling/sensortower\_state\_of\_ai\_2026/})에 저장된다. CLAUDE.md의 ``가공 데이터는 \texttt{data/proc/}'' 규칙의 예외이며 manifest에 기록했다.

% ============================================================
\section{2단계: 메시지 수(\texttt{03\_merge\_sensortower\_openai\_messages.do})}
% ============================================================

\subsection{가정}
\begin{description}[leftmargin=1.2cm,style=nextline]
  \item[(A1)] \textbf{모든 제품의 이용자 1인당 월 메시지 수는 같다.} 구현상 이 값은 시점과 국가에 따라서도 변하지 않는다. 2025년 7월 ChatGPT 값 하나를 2023-04--2026-05의 모든 제품$\cdot$국가에 쓴다.
  \item[(A2)] \textbf{모든 제품의 업무용 메시지 비중은 같다.} 그 비중은 국가별로 Signals의 ChatGPT 업무 비중과 같다고 둔다.
\end{description}

\subsection{이용자 1인당 메시지 수}
\begin{table}[h]
\centering
\caption{이용자 1인당 월 메시지 수 산출}\label{tab:anchor}
\begin{tabular}{lr}
\toprule
항목 & 값 \\
\midrule
% 04_analysis_sensortower_messages_summary.do가 생성. 직접 고치지 말 것.
ChatGPT 주 메시지(억 건, 2025-07, 전 세계) & 180.0 \\
ChatGPT 월 메시지(억 건, 전 세계) = 주 메시지 $\times$ 31/7 & 797.1 \\
18개국 비중(\%, AEI Claude.ai 2025-08-04--08-11) & 61.01 \\
ChatGPT 월 메시지(억 건, 18개국) = 전 세계 $\times$ 18개국 비중 & 486.3 \\
ChatGPT True Audience 이용자(억 명, 18개국 합계) & 9.05 \\
이용자 1인당 월 메시지(건) & 53.75 \\

\bottomrule
\end{tabular}
\end{table}

2025년 7월 월 메시지는 주 180억 건 $\times$ 31/7이다. 분모는 ⑯의 ChatGPT 2025년 7월 Worldwide True Audience 이용자다. 즉
\[
  m \equiv \text{msgs\_per\_user} = \frac{180\text{억} \times 31/7}{\text{users}_{\text{WW},\text{ChatGPT},2025\text{-}07}} .
\]

\subsection{제품별$\cdot$국가별$\cdot$Worldwide 메시지 수}
\[
  \text{messages}_{p,t} = m \cdot \text{users}_{\text{WW},p,t}, \quad
  \text{messages}_{c,t} = m \cdot \text{all\_users}_{c,t}, \quad
  \text{messages}_{\text{WW},t} = m \cdot \textstyle\sum_p \text{users}_{\text{WW},p,t} .
\]
국가별 값은 residual 행을 포함한다. do 파일은 매월 ``제품 합계 = 국가 합계(18개국 + residual) = Worldwide''를 검사한다. 2025년 7월 ChatGPT 메시지는 기준값 797.1억 건과 같다(검사).

\begin{table}[h]
\centering
\caption{제품별 Worldwide 이용자와 월간 메시지}\label{tab:products}
\begin{tabular}{lrrrrrr}
\toprule
 & \multicolumn{3}{c}{이용자(억 명)} & \multicolumn{3}{c}{월간 메시지(억 건)} \\
\cmidrule(lr){2-4}\cmidrule(lr){5-7}
제품 & 2024-07 & 2025-07 & 2026-05 & 2024-07 & 2025-07 & 2026-05 \\
\midrule
% 04_analysis_sensortower_messages_summary.do가 생성. 직접 고치지 말 것.
ChatGPT & 4.62 & 9.05 & 10.11 & 248.09 & 486.31 & 543.32 \\
Google Gemini & 0.79 & 3.08 & 6.05 & 42.73 & 165.74 & 324.99 \\
Claude & 0.23 & 0.45 & 2.22 & 12.10 & 24.08 & 119.39 \\
Microsoft Copilot & 0.21 & 0.37 & 0.36 & 11.09 & 19.84 & 19.12 \\
Perplexity & 0.20 & 0.93 & 0.62 & 11.01 & 49.96 & 33.12 \\
DeepSeek & 0.00 & 0.59 & 0.47 & 0.12 & 31.84 & 25.08 \\
Grok & 0.01 & 0.54 & 0.68 & 0.64 & 29.07 & 36.66 \\
Other & 0.41 & 0.54 & 1.04 & 22.04 & 29.24 & 56.01 \\
\midrule
합계 & 6.47 & 15.56 & 21.54 & 347.82 & 836.08 & 1157.69 \\

\bottomrule
\end{tabular}
\end{table}

(A1) 때문에 제품별 메시지 비중은 이용자 비중(⑯)과 똑같다. 예를 들어 Claude의 2026년 5월 메시지 비중은 이용자 비중 10.3\%와 같다. 제품 간 1인당 사용 강도 차이(report 15에서 본 앱 세션$\cdot$웹 방문 차이 등)는 반영되지 않는다.

\subsection{업무용 메시지 수}
\[
  \text{work\_messages}_{c,t} = \text{messages}_{c,t} \times \text{work\_share}_{c,t} .
\]
$\text{work\_share}_{c,t}$는 Signals 국가별 \texttt{work\_related} = 1 비중이다. residual 행(나머지 7개 시장)은 국가를 알 수 없어 Signals \textbf{전 세계} 비중을 쓴다(\texttt{work\_share\_src} = 2). Worldwide 업무용 메시지는 18개국과 residual의 합이며, 비교용으로 ``Worldwide 메시지 $\times$ Signals 전 세계 비중''(\texttt{work\_messages\_global})도 저장한다. Signals가 2024-07부터 있으므로 업무용 파일은 2024-07--2026-05(23개월)만 포함한다.

매칭은 Sensor Tower 시장명을 ISO 3166-1 alpha-2로 바꾼 코드와 월(Signals의 \texttt{YYYY-MM-01}에서 앞 7자리)로 한다. 18개국 모두 23개월 내내 Signals 값이 있다(do 파일에서 검사).

\begin{table}[h]
\centering
\small
\caption{국가별 월간 메시지와 업무용 메시지(2026년 5월)}\label{tab:country}
\begin{tabular}{lrrrrrr}
\toprule
시장 & 이용자 & 점유율 & 메시지 & 업무 비중 & 업무용 메시지 & 업무용 대화 \\
 & (백만 명) & (\%) & (억 건) & (\%) & (억 건) & ($k=3$, 억 건) \\
\midrule
% 04_analysis_sensortower_messages_summary.do가 생성. 직접 고치지 말 것.
India & 724.0 & 30.3 & 389.2 & 33.2 & 129.2 & 43.1 \\
United States & 293.6 & 12.3 & 157.8 & 32.1 & 50.7 & 16.9 \\
Brazil & 209.6 & 8.8 & 112.6 & 35.4 & 39.9 & 13.3 \\
Indonesia & 126.9 & 5.3 & 68.2 & 40.8 & 27.8 & 9.3 \\
Japan & 101.7 & 4.3 & 54.7 & 27.8 & 15.2 & 5.1 \\
Mexico & 92.6 & 3.9 & 49.8 & 29.1 & 14.5 & 4.8 \\
Turkey & 74.9 & 3.1 & 40.3 & 18.6 & 7.5 & 2.5 \\
Germany & 69.1 & 2.9 & 37.1 & 26.1 & 9.7 & 3.2 \\
South Korea & 61.0 & 2.6 & 32.8 & 31.7 & 10.4 & 3.5 \\
France & 59.8 & 2.5 & 32.1 & 25.5 & 8.2 & 2.7 \\
Vietnam & 58.3 & 2.4 & 31.4 & 31.8 & 10.0 & 3.3 \\
United Kingdom & 56.3 & 2.4 & 30.3 & 32.6 & 9.9 & 3.3 \\
Spain & 53.3 & 2.2 & 28.7 & 29.8 & 8.5 & 2.8 \\
Italy & 47.6 & 2.0 & 25.6 & 30.5 & 7.8 & 2.6 \\
Canada & 39.8 & 1.7 & 21.4 & 32.9 & 7.0 & 2.3 \\
Argentina & 38.4 & 1.6 & 20.7 & 29.1 & 6.0 & 2.0 \\
Australia & 26.1 & 1.1 & 14.0 & 39.7 & 5.6 & 1.9 \\
Taiwan & 20.7 & 0.9 & 11.1 & 33.5 & 3.7 & 1.2 \\
\midrule
나머지 7개 시장(residual) & 235.3 & 9.8 & 126.4 & 30.5 & 38.6 & 12.9 \\
\midrule
Worldwide(합계) & 2389.1 & 100.0 & 1284.1 & 31.9 & 410.1 & 136.7 \\

\bottomrule
\end{tabular}
\par\smallskip\footnotesize 주: 이용자는 8개 제품 True Audience 이용자 수의 합(중복 계산 포함). 업무 비중은 Signals 국가별 ChatGPT 값이며, residual은 Signals 전 세계 값, Worldwide 행은 함의된 비중(업무용 메시지 $\div$ 메시지).
\end{table}

\subsection{업무용 대화량}
대화당 메시지 수 $k$에 대한 공개 자료가 없으므로 격자로 둔다.
\[
  \text{conv\_k}K_{c,t} = \text{work\_messages}_{c,t} / K, \qquad K \in \{1.5, 3, 5, 7, 10\}
\]
변수 이름은 \texttt{conv\_k1p5}, \texttt{conv\_k3}, \texttt{conv\_k5}, \texttt{conv\_k7}, \texttt{conv\_k10}이다. 국가별$\cdot$Worldwide 업무용 파일에 모두 들어 있다.

\begin{table}[h]
\centering
\footnotesize
\caption{Worldwide 월간 메시지, 업무용 메시지, 업무용 대화}\label{tab:ww}
\begin{tabular}{lrrrrrrrrr}
\toprule
 & 메시지 & \multicolumn{2}{c}{업무 비중(\%)} & 업무용 메시지 & \multicolumn{5}{c}{업무용 대화(억 건), 대화당 메시지 $k$} \\
\cmidrule(lr){3-4}\cmidrule(lr){6-10}
월 & (억 건) & 함의 & Signals 전 세계 & (억 건) & 1.5 & 3 & 5 & 7 & 10 \\
\midrule
% 04_analysis_sensortower_messages_summary.do가 생성. 직접 고치지 말 것.
2024-07 & 566.6 & 53.8 & 51.3 & 305.1 & 203.4 & 101.7 & 61.0 & 43.6 & 30.5 \\
2024-08 & 605.8 & 51.1 & 48.3 & 309.7 & 206.5 & 103.2 & 61.9 & 44.2 & 31.0 \\
2024-09 & 684.5 & 49.3 & 46.6 & 337.8 & 225.2 & 112.6 & 67.6 & 48.3 & 33.8 \\
2024-10 & 749.2 & 48.1 & 45.4 & 360.5 & 240.3 & 120.2 & 72.1 & 51.5 & 36.0 \\
2024-11 & 805.5 & 46.8 & 43.7 & 377.2 & 251.5 & 125.7 & 75.4 & 53.9 & 37.7 \\
2024-12 & 828.1 & 43.2 & 40.1 & 357.5 & 238.3 & 119.2 & 71.5 & 51.1 & 35.7 \\
2025-01 & 940.4 & 43.6 & 40.3 & 409.9 & 273.3 & 136.6 & 82.0 & 58.6 & 41.0 \\
2025-02 & 1050.1 & 41.9 & 38.9 & 440.1 & 293.4 & 146.7 & 88.0 & 62.9 & 44.0 \\
2025-03 & 1137.7 & 40.0 & 37.6 & 455.3 & 303.5 & 151.8 & 91.1 & 65.0 & 45.5 \\
2025-04 & 1203.6 & 37.8 & 35.2 & 454.6 & 303.1 & 151.5 & 90.9 & 64.9 & 45.5 \\
2025-05 & 1257.2 & 38.0 & 35.5 & 477.3 & 318.2 & 159.1 & 95.5 & 68.2 & 47.7 \\
2025-06 & 1293.2 & 35.7 & 33.3 & 461.6 & 307.8 & 153.9 & 92.3 & 65.9 & 46.2 \\
2025-07 & 1374.3 & 34.1 & 31.4 & 468.6 & 312.4 & 156.2 & 93.7 & 66.9 & 46.9 \\
2025-08 & 1437.7 & 32.9 & 30.6 & 472.3 & 314.9 & 157.4 & 94.5 & 67.5 & 47.2 \\
2025-09 & 1534.1 & 35.0 & 33.0 & 537.2 & 358.1 & 179.1 & 107.4 & 76.7 & 53.7 \\
2025-10 & 1578.3 & 35.8 & 33.9 & 564.3 & 376.2 & 188.1 & 112.9 & 80.6 & 56.4 \\
2025-11 & 1593.0 & 35.2 & 33.6 & 561.3 & 374.2 & 187.1 & 112.3 & 80.2 & 56.1 \\
2025-12 & 1596.9 & 32.4 & 30.5 & 517.9 & 345.2 & 172.6 & 103.6 & 74.0 & 51.8 \\
2026-01 & 1630.4 & 32.7 & 30.7 & 532.7 & 355.2 & 177.6 & 106.5 & 76.1 & 53.3 \\
2026-02 & 1644.9 & 32.5 & 30.8 & 533.8 & 355.9 & 177.9 & 106.8 & 76.3 & 53.4 \\
2026-03 & 1780.7 & 32.6 & 31.0 & 580.1 & 386.7 & 193.4 & 116.0 & 82.9 & 58.0 \\
2026-04 & 1866.1 & 33.3 & 31.7 & 621.0 & 414.0 & 207.0 & 124.2 & 88.7 & 62.1 \\
2026-05 & 1915.9 & 31.9 & 30.5 & 611.9 & 407.9 & 204.0 & 122.4 & 87.4 & 61.2 \\

\bottomrule
\end{tabular}
\par\smallskip 주: ``함의'' = 18개국 + residual 업무용 메시지 합 $\div$ Worldwide 메시지.
\end{table}

함의된 Worldwide 업무 비중은 모든 달에서 Signals 전 세계 비중보다 1.4--3.3\%p 높다. Sensor Tower 25개 시장의 이용자 구성이 Signals 전 세계 메시지 구성과 다르기 때문이다(예: 2026년 5월 이용자 비중 30\%인 인도의 업무 비중 33.2\%가 전 세계 30.5\%보다 높다). 업무 비중은 2024년 7월 약 54\%에서 2026년 5월 약 32\%로 낮아졌는데, 이는 Signals ChatGPT 업무 비중 하락을 그대로 옮긴 것이다.

% ============================================================
\section{산출 파일}
% ============================================================

\begin{longtable}{L{5.4cm}L{6.2cm}R{1.2cm}L{2.6cm}}
\caption{산출 파일(\texttt{data/proc/macro/scaling/sensortower\_state\_of\_ai\_2026/}, 각 .csv/.dta)}\label{tab:files}\\
\toprule
파일 & 주요 변수 & 행 수 & 키 \\
\midrule
\endfirsthead
\toprule
파일 & 주요 변수 & 행 수 & 키 \\
\midrule
\endhead
\texttt{sensortower\_messages\_by\_product\_monthly} & \texttt{unique\_users}, \texttt{ww\_total\_users}, \texttt{share\_ww\_pct}, \texttt{msgs\_per\_user}, \texttt{messages} & 304 & assistant $\times$ month \\
\texttt{sensortower\_messages\_by\_country\_monthly} & \texttt{iso2}, \texttt{residual}, \texttt{all\_users}, \texttt{share\_all\_products\_pct}, \texttt{messages} & 722 & market $\times$ month \\
\texttt{sensortower\_messages\_worldwide\_monthly} & \texttt{worldwide\_users}, \texttt{msgs\_per\_user}, \texttt{messages} & 38 & month \\
\texttt{sensortower\_work\_messages\_by\_country\_monthly} & \texttt{work\_share}, \texttt{work\_share\_src}, \texttt{work\_messages}, \texttt{nonwork\_messages}, \texttt{conv\_k*} & 437 & market $\times$ month \\
\texttt{sensortower\_work\_messages\_worldwide\_monthly} & \texttt{work\_share\_implied}, \texttt{work\_messages}, \texttt{work\_share\_global}, \texttt{work\_messages\_global}, \texttt{conv\_k*} & 23 & month \\
\bottomrule
\end{longtable}

기간은 메시지 파일이 2023-04--2026-05(38개월), 업무용 파일이 2024-07--2026-05(23개월)다. 국가별 파일은 18개국 + residual(19행/월)이다. 요약표는 \texttt{results/table/17\_sensortower\_messages.xlsx}에도 있다.

% ============================================================
\section{해석상 주의와 한계}
% ============================================================

\begin{enumerate}[leftmargin=*]
  \item \textbf{(A1)의 강도.} 제품$\cdot$국가$\cdot$시점에 걸쳐 1인당 메시지가 같다고 두었다. 실제로는 ChatGPT의 메시지가 2024년 6월--2025년 7월 5.9배 늘 때 이용자는 1.9배 늘어 1인당 사용이 크게 변했다(report 15). 따라서 2025년 7월에서 멀어질수록 오차가 커지고, 특히 2024년 이전 메시지는 과대, 2026년 메시지는 과소일 가능성이 크다. 제품별로도 Claude 이용자의 대부분이 웹 이용자라는 점 등 사용 강도 차이가 있다.
  \item \textbf{(A2)의 강도.} 업무 비중을 ChatGPT 값으로 통일했다. 업무 중심 제품(Claude, Copilot)의 업무용 메시지는 과소, 일반 소비자 중심 제품이 많은 국가의 업무용 메시지는 과대일 수 있다. 이 점은 report 18에서 AEI와 비교해 점검한다.
  \item \textbf{범위 불일치.} 기준값 180억 건은 전 세계 ChatGPT(소비자 요금제) 값이지만 분모는 Sensor Tower 25개 시장의 독립 앱+웹 이용자다. 25개 시장 밖 이용자의 메시지가 분자에 들어가므로 1인당 메시지가 다소 크다. API, 기업 요금제, 앱 내장 AI 사용은 빠진다.
  \item \textbf{중복 계산.} 국가 메시지는 제품별 이용자 수 합에 1인당 값을 곱한 것이므로 ``그 국가 모든 제품 메시지의 합''이다. 여러 어시스턴트를 쓰는 사람의 메시지는 제품 수만큼 더해진다. 이는 (A1) 아래에서 일관적이다.
  \item \textbf{대화량.} 대화당 메시지 수 $k$는 근거 자료가 없어 1.5--10의 격자로만 제시한다. 대화량은 $k$에 반비례하므로 메시지 수보다 훨씬 불확실하다.
  \item \textbf{반올림과 잡음.} Sensor Tower 이용자 수(유효숫자 3자리)와 Signals 비중(차등정보보호 잡음, 소수 셋째 자리)의 오차가 그대로 전파된다.
\end{enumerate}

% ============================================================
\section{재현}
% ============================================================

Stata 17에서 다음 순서로 실행한다(모든 do 파일은 \texttt{\$root}를 기준으로 경로를 잡는다).
\begin{enumerate}[leftmargin=*]
  \item \texttt{code/02\_clean\_sensortower\_share\_all\_products.do} --- ⑮, ⑯(원자료 폴더에 저장).
  \item \texttt{code/02\_clean\_sensortower\_share\_within\_product.do} --- 제품 내 국가 비중(report 18에서 사용).
  \item \texttt{code/03\_merge\_sensortower\_openai\_messages.do} --- 위 다섯 데이터.
  \item \texttt{code/04\_analysis\_sensortower\_messages\_summary.do} --- 본 문서의 표(\texttt{results/table/17\_*.tex}, \texttt{.xlsx}).
\end{enumerate}

\printbibliography[title={참고문헌}]

\end{document}
